\documentclass[11pt,a4paper]{article}

\usepackage[utf8]{inputenc}
\usepackage[T1]{fontenc}
\usepackage{amsmath,amssymb}
\usepackage{graphicx}
\usepackage{booktabs}
\usepackage{hyperref}
\usepackage[margin=0.9in]{geometry}
\usepackage{caption}
\usepackage{subcaption}
\usepackage{float}
\usepackage{xcolor}
% Fallback for long underscored filenames when seqsplit.sty is unavailable:
% insert an allowed break after every character of the argument.
\makeatletter
\providecommand{\seqsplit}[1]{\@tfor\@tempc:=#1\do{\@tempc\allowbreak}}
\makeatother
\sloppy

\newcommand{\etg}{$E_T^{\gamma,\mathrm{rec}}$}
\newcommand{\wetacog}{\texorpdfstring{$w_{\eta}^{\mathrm{CoG}}$}{w\_eta\textasciicircum CoG}}
\newcommand{\wphicog}{\texorpdfstring{$w_{\phi}^{\mathrm{CoG}}$}{w\_phi\textasciicircum CoG}}
\newcommand{\isoET}{E_T^{\mathrm{iso}}}
\newcommand{\pp}{$p$+$p$}
\newcommand{\epem}{\ensuremath{e^{+}e^{-}}}

\title{Converted vs.\ Unconverted Photons: Shower Shape,\\
Isolation, and Energy Response in photon10 MC}
\author{PPG12 Analysis Team (automated report)}
\date{April 16, 2026}

\begin{document}
\maketitle

\begin{abstract}
We study the impact of photon conversions (pre-EMCal \epem\ pair production in
detector material) on reconstruction observables used by the PPG12 isolated-photon
cross-section measurement in \pp\ at $\sqrt{s}=200$~GeV.  Using the
\texttt{photon10} single-particle Monte Carlo ($\sim 10$\,M events, truth
\(p_T \in [14,30]\)~GeV), we find that \textbf{17.3\%} of fiducial truth
photons ($|\eta|<0.7$, \(14 \le p_T < 30\)~GeV) convert before the EMCal.
% source: fractions_findings.md Sec. inclusive |eta|<0.7
Converted photons exhibit (i) a mean cluster-energy response shift of
\(\Delta\langle R_E \rangle \simeq -0.076\) (\(-8\%\)) driven by an asymmetric
low-$R$ tail,
% source: response_findings.md Sec.4 and Sec.9 summary table
(ii) a drop of $-14.6$~percentage points in the parametric reco-isolation pass
fraction (rising to $-22.8$~pp at 8--10~GeV), dominated by additional EMCal
energy inside the isolation cone,
% source: isolation_findings.md Sec. parametric iso pass fraction
and (iii) a $-16\%$ shift in the mean shower-shape BDT score, with
$w_{\phi}^{\mathrm{CoG}}$ broadening by $+91\%$ due to the solenoidal-field
opening of the \epem\ pair.
% source: showershape_findings.md Sec.3 inclusive table, Sec.5 discriminating ranking
These four features are mutually consistent with a single physical picture:
the $\sim 1.4$~T sPHENIX solenoid splits the \epem\ pair in azimuth, one arm
leaks outside the cluster footprint (reducing the energy response and
polluting the isolation cone), and the resulting wider, two-lobed
$\phi$-profile degrades the ID BDT. The effect enters the cross-section
measurement as a joint ID$+$iso efficiency systematic whose size is set by
the accuracy of the simulated material budget.
\end{abstract}


%% ===============================================================
\section{Introduction and Motivation}
\label{sec:intro}

Photon conversion --- the reaction \(\gamma \to \epem\) induced by the Coulomb
field of a nucleus --- occurs with probability \(P \simeq 1 - \exp\left[-\tfrac{7}{9}\,X/X_0\right]\)
as a photon traverses a thickness \(X/X_0\) of material. In sPHENIX the tracker
(MVTX + INTT + TPC) together with its support structures and cables precedes the
barrel EMCal. With the nominally quoted $\sim 10$--$20$\%~$X_0$ of material, a
\(\sim 10\)--$15\%$ conversion probability is expected in the central barrel.
Once a photon converts, the two leptons are bent in opposite azimuthal
directions by the $\sim 1.4$~T solenoid, land in the EMCal as two partially
overlapping electromagnetic showers, and are reconstructed by the standard
clustering as a single (wider) cluster whose properties differ from those of
a pristine photon.

These distortions matter for PPG12 because the tight-selected isolated photon
signal efficiency is the product of a shower-shape identification efficiency
(via the $E_T$-dependent BDT cut) and an isolation efficiency (via the
parametric isolation cone cut). Both steps are sensitive to the converted
fraction and its pre-reconstruction kinematics. Any mismatch in the simulated
material budget between Monte Carlo and data therefore propagates directly
into the cross-section. Quantifying the conversion fingerprint on each
observable --- and its \(p_T\) and \(\eta\) dependence --- is a prerequisite
for assigning a material-budget systematic uncertainty.

Truth conversion information is carried in the slimtree through the
\texttt{particle\_converted} flag written by
\texttt{anatreemaker/source/CaloAna24.cc}.  The flag is built by traversing the
G4 shower particle list of each truth photon
(\texttt{CaloAna24.cc} lines \(\sim 807\text{--}843\)): any secondary that
(a) originates at a G4 vertex with $r_\perp<93$~cm (i.e.\ upstream of the EMCal
front face) and (b) carries more than 40\% of the photon momentum, tags the
photon.  If at least one such secondary is an \(e^{\pm}\)
(\(|\mathrm{pid}|=11\)), the photon is labelled \texttt{converted}; otherwise
it is labelled \texttt{bad}.  The flag is written at
\texttt{CaloAna24.cc} lines \(\sim 953\text{--}970\) using:
%
\begin{equation}
\texttt{particle\_converted} =
\begin{cases}
0 & \text{no upstream secondary above threshold (unconverted)} \\
1 & \text{\epem\ daughter above threshold (converted)} \\
2 & \text{non-\epem\ secondary above threshold (bad)}
\end{cases}
\end{equation}
%
In what follows we treat the small ``bad'' population (0.01--0.02\% of truth
photons, $<300$ matched in 10~M events) as physically negligible and focus on
the unconverted--converted comparison.


%% ===============================================================
\section{Data Sample and Method}
\label{sec:sample}

\paragraph{Sample.} The study is based on the single-photon Monte Carlo file
\texttt{\seqsplit{photon10\_with\_bdt\_vtx\_noreweight\_single.root}} (one
fired photon per event, truth $p_T \in [14,30]$~GeV at generation,
$|\eta|<1.5$).
Four independent analyses read the tree \texttt{slimtree}:
% source: showershape_findings.md header; isolation_findings.md header; response_findings.md header; fractions_findings.md header
\begin{itemize}
  \item Fractions, shower shape, and energy response use the full 9{,}998{,}541
    events.
  \item Isolation uses a 2{,}000{,}000-event subsample (runtime-based).
\end{itemize}

\paragraph{Truth selection (fiducial).} $\texttt{particle\_pid} == 22$,
$|\texttt{particle\_Eta}| < 0.7$, $\texttt{particle\_truth\_iso\_03} < 4$~GeV,
and (where applicable) $8 \le \texttt{particle\_Pt} \le 36$~GeV or
$14 \le \texttt{particle\_Pt} < 30$~GeV as indicated.
% source: showershape_findings.md Sec.1; fractions_findings.md Sec. inclusive |eta|<0.7

\paragraph{Cluster--truth matching.} For each truth photon the highest-$E_T$
cluster in \texttt{CLUSTERINFO\_CEMC} whose
\texttt{cluster\_truthtrkID} equals the truth \texttt{particle\_trkid} is
selected. A residual cluster $E_T$ threshold ($>3$--$5$~GeV, depending on the
analysis) is applied. pT binning follows the PPG12 analysis edges
\(\texttt{ptRanges}=\{8,10,12,14,16,18,20,22,24,26,28,32,36\}\) (12 bins) with
the matched cluster $E_T$ as the abscissa.
% source: showershape_findings.md Sec.1; response_findings.md header

\paragraph{Caveat on isolation.} The input file \emph{does not} contain the
topo-cluster isolation branches
\texttt{cluster\_iso\_topo\_03/04}. The isolation study therefore uses the
calo-tower isolation \texttt{cluster\_iso\_03} (and R=0.4 equivalent). The
nominal analysis configuration \texttt{efficiencytool/config\_bdt\_nom.yaml}
uses topo-iso at $R=0.4$ (\texttt{use\_topo\_iso: 2}). The \emph{shape} of the
converted/unconverted difference is expected to carry over, but absolute
magnitudes should be re-measured on topo-iso branches once they are populated
for \texttt{photon10}.
% source: isolation_findings.md Selection summary, Concerns and next steps

\paragraph{Shower shape caveat.} The shower-shape BDT score is only written to
the \texttt{CLUSTERINFO\_CEMC\_NO\_SPLIT} node, whereas the other shower-shape
observables use \texttt{CLUSTERINFO\_CEMC}. BDT conclusions below are therefore
based on the unsplit-cluster node.
% source: showershape_findings.md Sec. header


%% ===============================================================
\section{Conversion Fractions and Cluster Match Efficiency}
\label{sec:fractions}

In the fiducial region ($|\eta|<0.7$, $14 \le p_T < 30$~GeV, 855{,}380 truth
photons), the converted fraction is $(17.317 \pm 0.041)\%$ and the ``bad''
fraction is a negligible $(0.024 \pm 0.002)\%$.
% source: fractions_findings.md Sec. "Inclusive (|eta|<0.7, 14<=pT<30 GeV)"
Table~\ref{tab:fractions_inclusive} summarises the inclusive category counts
and per-category cluster-match efficiency. The converted fraction is
effectively flat in \(p_T\) ($17.2$--$17.9\%$ across 8--36~GeV) and rises with
$|\eta|$, from $(15.4 \pm 0.3)\%$ at $|\eta|<0.1$ to $(20.1 \pm 0.4)\%$ at
$0.6 < |\eta| < 0.7$, reaching $32\%$ at $|\eta| \simeq 1.4$--$1.5$
(extended window).
% source: fractions_findings.md Sec. per-|eta|-bin tables

\begin{table}[H]
\centering
\caption{Converted-photon fractions and cluster-match efficiency in the
fiducial window $|\eta|<0.7$, $14 \le p_T < 30$~GeV.  Match efficiency is
defined as the fraction of truth photons (within a category) for which at
least one cluster satisfies $\texttt{cluster\_truthtrkID}==\texttt{particle\_trkid}$
and $\texttt{cluster\_Et}>5$~GeV.
% source: fractions_findings.md "Inclusive (|eta|<0.7, 14<=pT<30 GeV)" and "Cluster-match efficiency"
}
\label{tab:fractions_inclusive}
\begin{tabular}{lrrr}
\toprule
Category & N (truth) & Fraction & Match efficiency \\
\midrule
Unconverted         & 707{,}047 & $(82.659 \pm 0.041)\%$ & $(67.987 \pm 0.055)\%$ \\
Converted (\epem)   & 148{,}129 & $(17.317 \pm 0.041)\%$ & $(66.290 \pm 0.123)\%$ \\
Bad secondary       &       204 & $(0.024  \pm 0.002)\%$ & $(39.024 \pm 3.407)\%$ \\
\bottomrule
\end{tabular}
\end{table}

\noindent The three panels of Fig.~\ref{fig:fractions_panels} show the main
trends:
\begin{figure}[H]
\centering
\begin{subfigure}[t]{0.32\textwidth}
\includegraphics[width=\textwidth]{converted_photon_study/figures/fractions_vs_pt_eta07.pdf}
\caption{Fractions vs.\ \(p_T\) ($|\eta|<0.7$)}
\end{subfigure}
\hfill
\begin{subfigure}[t]{0.32\textwidth}
\includegraphics[width=\textwidth]{converted_photon_study/figures/fractions_vs_abseta_eta07.pdf}
\caption{Fractions vs.\ $|\eta|$ ($14 \le p_T < 30$~GeV)}
\end{subfigure}
\hfill
\begin{subfigure}[t]{0.32\textwidth}
\includegraphics[width=\textwidth]{converted_photon_study/figures/fractions_match_eff_eta07.pdf}
\caption{Cluster-match efficiency vs.\ \(p_T\)}
\end{subfigure}
\caption{Converted-photon fractions and match efficiency in the fiducial
window. The converted fraction is flat in \(p_T\) and rises with $|\eta|$
consistent with increased integrated material at forward rapidity. Converted
photons have $\sim 2$~pp lower cluster-match efficiency than unconverted.
% source: fractions_findings.md Sec. "Per-pT-bin fractions" and "Per-|eta|-bin fractions"
}
\label{fig:fractions_panels}
\end{figure}

The central-barrel converted fraction of $\sim 17\%$ is consistent with a
material budget near the upper end of the nominally quoted $10$--$20\%\,X_0$
range ($1-\exp[-\tfrac{7}{9}\cdot 0.20] \simeq 15\%$). The sharp rise toward
$|\eta|\sim 1.5$ is consistent with endcap-style material growth and matches
expectations for the tracker support geometry.
% source: fractions_findings.md Sec. "Physics discussion"

The absolute match-efficiency value ($\sim 66$--$68\%$ for fiducial photons)
looks low at first glance because the definition \emph{requires}
\texttt{cluster\_truthtrkID} to equal the \emph{primary} photon track id. When
clustering associates a post-shower secondary ($e^\pm$, brem $\gamma$) instead
of the original photon, the match fails even though a physically correct
cluster is present. The converted-category match efficiency is therefore a
\emph{lower bound}; the \emph{relative} $\sim 2.5\%$ deficit with respect to
unconverted is the physics-meaningful quantity.
% source: fractions_findings.md Sec. "Concerns / caveats"


%% ===============================================================
\section{Energy Response and Resolution}
\label{sec:response}

We define the cluster energy response \(R_E = E_{\mathrm{cluster}}/E_{\gamma}^{\mathrm{truth}}\)
and compare its shape, first and second moments, core Gaussian and Crystal
Ball fit parameters, and low-$R$ tail fractions between the converted and
unconverted categories. Analogous results for $R_{E_T}$ are contained in the
auxiliary grid \texttt{response\_shape\_ET\_grid.pdf}.
% source: response_findings.md Sec. "Energy response shapes"
The abscissa is the \emph{truth} photon $p_T$ (so that the per-bin
efficiency reflects the pre-clustering kinematics); the isolation and
shower-shape sections that follow bin in the matched cluster $E_T$ instead,
which is the variable observed in data.
The negligible ``bad'' category (see Sec.~\ref{sec:intro}) is suppressed from
the summary plots and from the Crystal Ball fits but remains in the shape
grids for completeness.

\subsection{Shape, mean, and resolution}

\begin{figure}[H]
\centering
\begin{subfigure}[t]{0.48\textwidth}
\includegraphics[width=\textwidth]{converted_photon_study/figures/response_shape_grid.pdf}
\caption{Response shape \(R_E\) per \(p_T\) bin}
\end{subfigure}
\hfill
\begin{subfigure}[t]{0.48\textwidth}
\includegraphics[width=\textwidth]{converted_photon_study/figures/response_mean_vs_pt.pdf}
\caption{Mean response \(\langle R_E \rangle\) vs.\ \(p_T\)}
\end{subfigure}
\caption{Response distributions and means for unconverted vs.\ converted
photons. The converted distribution shows a pronounced low-$R$ tail at
$p_T \lesssim 20$~GeV caused by one arm of the \epem\ pair leaking outside
the cluster footprint. Above $p_T\sim 25$~GeV the distributions converge and
the mean response of converted photons crosses above that of unconverted
(pair-collimation limit).
% source: response_findings.md Sec. 3 and Sec. 4
}
\label{fig:response_shape_and_mean}
\end{figure}

\begin{figure}[H]
\centering
\begin{subfigure}[t]{0.48\textwidth}
\includegraphics[width=\textwidth]{converted_photon_study/figures/response_sigma_vs_pt.pdf}
\caption{Gaussian and Crystal Ball core $\sigma$ vs.\ \(p_T\)}
\end{subfigure}
\hfill
\begin{subfigure}[t]{0.48\textwidth}
\includegraphics[width=\textwidth]{converted_photon_study/figures/response_tailfrac_vs_pt.pdf}
\caption{Tail fractions $R<0.8$ and $R<0.5$ vs.\ \(p_T\)}
\end{subfigure}
\caption{Resolution and tail-fraction summary. The Gaussian core width
(fit in $[\mathrm{peak}-0.10,\,\mathrm{peak}+0.10]$) is only a few percent
wider for converted photons, reflecting the fact that the surviving \epem\
core reconstructs at essentially the same energy as an unconverted photon.
The Crystal Ball $\sigma$ (simultaneous fit of Gaussian core plus left-side
power-law tail over $R\in[0.3,1.2]$) shows a clear broadening of the core
at low \(p_T\) --- up to $\sim +70\%$ at $p_T \sim 11$~GeV --- that converges
to the unconverted value above $\sim 20$~GeV. The Gaussian-$\sigma$ panel
shows a mild excursion for the converted category at $p_T \sim 22$--$32$~GeV:
the symmetric $\pm 0.10$ window is fragile in that stats-limited range
because the peak slides within the window as the converted sub-peak
structure shifts; the Crystal Ball $\sigma$, which fits the core independently
of the tail, is the stable estimator there. Both fits exclude the
``bad'' category.
% source: summary.txt, response_cb_params_vs_pt.pdf
}
\label{fig:response_sigma_tail}
\end{figure}

\begin{figure}[H]
\centering
\includegraphics[width=0.85\textwidth]{converted_photon_study/figures/response_cb_params_vs_pt.pdf}
\caption{Crystal Ball tail-transition parameter $\alpha$ (in units of the
fitted $\sigma$) and power-law index $n$ vs.\ \(p_T\). The tail transition
$\alpha \simeq 0.45$--$0.6\,\sigma$ is similar between categories --- the
tail sets in very close to the peak in both --- while the core width itself
is what differs (Fig.~\ref{fig:response_sigma_tail}). The power-law index
$n$ is poorly constrained and typically hits the numerical bound; the
fit is effectively Gaussian-equivalent in the tail region for the
inclusive \(p_T\) distribution, so the physics information resides in
$\mu$, $\sigma$, and $\alpha$.
% source: response_cb_params_vs_pt.pdf, summary.txt
}
\label{fig:response_cb_params}
\end{figure}

\begin{table}[H]
\centering
\caption{Inclusive ($10 \le p_T \le 34$~GeV) energy-response summary for
converted vs.\ unconverted photons. The Gaussian fit is performed in a
symmetric window $[\mathrm{peak}-0.10,\,\mathrm{peak}+0.10]$ around the
3-bin-smoothed peak (core-only); the Crystal Ball (CB) fit uses the full
range $R\in[0.3,1.2]$ and parametrises a Gaussian core plus a left-side
power-law tail with transition at $\mu-\alpha\,\sigma$. The CB power-law
index $n$ is not well-constrained in the inclusive distribution (bumps
against the numerical upper bound); the physical information is carried by
$(\mu,\sigma,\alpha)$.
% source: summary.txt, summary.pkl (Gaussian + Crystal Ball fits in make_plots.py)
}
\label{tab:response_inclusive}
\begin{tabular}{lrrr}
\toprule
Quantity & Unconverted & Converted & Converted $-$ Unconverted \\
\midrule
$\langle R_E \rangle$               & 0.9424 & 0.8666 & $-0.0758$ ($-8\%$) \\
median $R_E$                         & 0.9650 & 0.9050 & $-0.060$ \\
upper width $p_{84}-p_{50}$         & 0.080  & 0.120  & $+50\%$ \\
lower width $p_{50}-p_{16}$         & 0.120  & 0.230  & $+92\%$ \\
Gaussian core $\mu$                  & 0.9894 & 0.9924 & $+0.003$ \\
Gaussian core $\sigma$               & 0.0653 & 0.0677 & $+4\%$ \\
Crystal Ball $\mu$                   & 1.0025 & 0.9741 & $-0.028$ \\
Crystal Ball $\sigma$                & 0.0514 & 0.0696 & $+35\%$ \\
Crystal Ball $\alpha$                & 0.463  & 0.458  & $-0.01$ \\
fraction $R<0.8$                    & 0.107  & 0.311  & $\times 2.9$ \\
fraction $R<0.5$                    & 0.0053 & 0.0201 & $\times 3.8$ \\
\bottomrule
\end{tabular}
\end{table}

The mean shift of $\Delta\langle R_E \rangle = -0.076$ is not a shift of the
Gaussian core: the Gaussian-core $\mu$ (symmetric $\pm0.10$ fit) moves by
only $+0.003$ and the Gaussian-core $\sigma$ broadens by $\sim 4\%$. The
Crystal Ball fit, which accommodates the left-side power-law tail, finds a
core $\mu$ shift of $-0.028$ and a core $\sigma$ broadening of $+35\%$ ---
consistent with the picture that a non-trivial fraction of the converted
population builds a sub-peak below $R=1$ while the surviving two-arm core
still sits close to unity. Quantitatively, the tail fraction below $R=0.8$
is $\sim 3\times$ larger for converted photons and below $R=0.5$ is
$\sim 4\times$ larger; these tails dominate the $-8\%$ mean shift.
% source: summary.txt (Gaussian + Crystal Ball fit parameters)

\subsection{\(p_T\) trend and crossover}

Per-$p_T$ shifts (from \texttt{response\_findings.md} Sec.~4):
$\Delta\langle R_E \rangle \simeq -0.082$ at 10--12~GeV, shrinking to
$-0.042$ at 18--20~GeV, $-0.014$ at 22--24~GeV, and \emph{crossing sign} at
$\sim 26$--$27$~GeV. Above $p_T \simeq 25$~GeV the mean response of converted
photons \emph{exceeds} that of unconverted (by up to $+0.023$ at 28--32~GeV).
Physically, the \epem\ opening angle after magnetic-field bending decreases
as $\sim 1/p_T$, so the pair becomes collimated enough to be contained in a
single cluster --- exactly where the \emph{unconverted} response continues to
degrade slowly with \(p_T\) from shower leakage into the HCal. The same
crossover appears in the tail-0.8 curves (converted tail decreases from 0.32
to 0.13 as \(p_T\) rises; unconverted tail rises mildly from 0.10 to 0.17),
confirming the effect.
% source: response_findings.md Sec. 4 and Sec. 6

\subsection{Angular deltas}

\begin{figure}[H]
\centering
\includegraphics[width=0.72\textwidth]{converted_photon_study/figures/response_angular_deltas.pdf}
\caption{Cluster-minus-truth angular deltas $\Delta\eta$, $\Delta\phi$ and
$\Delta R$. The $\Delta\phi$ distribution for converted photons is strikingly
bimodal with lobes at $\Delta\phi \sim \pm 0.01$~rad and a dip at zero ---
the signature of magnetic-field bending of the \epem\ pair, with the cluster
centroid shifting toward the higher-energy lepton.  $\Delta\eta$ shows no
shift or broadening, consistent with a purely azimuthal magnetic field.
$\Delta R$ has a broader tail for converted photons (up to $0.15$~rad) that
correlates with the $\Delta\phi$ lobes.
% source: response_findings.md Sec. 8
}
\label{fig:response_angular}
\end{figure}

The $\Delta\phi$ bimodality is a purely kinematic consequence of the
\epem\ pair carrying opposite charges: the solenoidal field bends the
leptons in opposite senses with equal probability, so the mean $\Delta\phi$
is zero while the RMS is roughly an order of magnitude larger than that of
unconverted photons.  This feature is used below (Sec.~\ref{sec:showershape})
to interpret the \wphicog\ distributions.


%% ===============================================================
\section{Isolation Energy}
\label{sec:isolation}

The PPG12 nominal reconstruction isolation cut is the parametric form
\begin{equation}
  \isoET\text{ max} = 0.502095 + 0.0433036 \cdot E_T,
  \label{eq:iso_cut}
\end{equation}
% source: efficiencytool/config_bdt_nom.yaml lines 33-34 (reco_iso_max_b, reco_iso_max_s)
given in GeV with $E_T$ also in GeV.  In the nominal analysis the input
variable is the topo-cluster isolation
\texttt{cluster\_iso\_topo\_04} (\texttt{use\_topo\_iso: 2} in
\texttt{config\_bdt\_nom.yaml}).
% source: efficiencytool/config_bdt_nom.yaml line 53
Because the input MC file lacks that branch, we apply the same cut to
\texttt{cluster\_iso\_03} (calo-tower iso, own-cluster towers removed).
Absolute numbers below are tied to this substitution; the qualitative
conclusions on converted--unconverted separation are robust because both
isolation definitions sum tower $E_T$ in an $R=0.3$ cone.

\subsection{Parametric cut pass fraction}

\begin{figure}[H]
\centering
\begin{subfigure}[t]{0.48\textwidth}
\includegraphics[width=\textwidth]{converted_photon_study/figures/isolation_cluster_iso_03.pdf}
\caption{\(\texttt{cluster\_iso\_03}\) shape per \(p_T\) bin}
\end{subfigure}
\hfill
\begin{subfigure}[t]{0.48\textwidth}
\includegraphics[width=\textwidth]{converted_photon_study/figures/isolation_pass_fraction_vs_pt.pdf}
\caption{Parametric-cut pass fraction vs.\ \(p_T\)}
\end{subfigure}
\caption{Isolation distributions (left) and parametric-cut pass fraction
(right). The low-\(p_T\) converted population has a long tail to $\isoET>2$~GeV
that fails the parametric cut; the inclusive pass-fraction deficit is
$-14.6$~pp for converted photons, dominated by the 8--14 GeV region.
% source: isolation_findings.md Sec. "Parametric iso pass fraction per pT bin"
}
\label{fig:iso_shape_pass}
\end{figure}

Inclusive ($8 \le E_T < 36$~GeV) pass fractions:
$\texttt{unconv}=0.649$, $\texttt{conv}=0.503$, $\Delta=-0.146$ (i.e.\ a
\textbf{$-14.6$~percentage-point} drop).
% source: isolation_findings.md Sec. "Parametric iso pass fraction per pT bin" inclusive row
The deficit peaks at $-22.8$~pp in the 8--10~GeV bin ($0.562\to 0.334$) and
shrinks to the $\pm 5$~pp statistics-limited level above 24~GeV.

\subsection{Mean isolation energy and calorimeter decomposition}

\begin{figure}[H]
\centering
\includegraphics[width=0.75\textwidth]{converted_photon_study/figures/isolation_mean_vs_pt_R003.pdf}
\caption{Mean isolation energy $\langle\texttt{cluster\_iso\_03}\rangle$
vs.\ cluster $E_T$, split into total + EMCal + inner HCal + outer HCal
sub-components.  The total and EMCal sub-components exhibit a clear
converted/unconverted separation; both HCal sub-components show the opposite
(small) trend.
% source: isolation_findings.md Sec. "Which calorimeter drives the shift?"
}
\label{fig:iso_mean_vs_pt}
\end{figure}

\begin{table}[H]
\centering
\caption{Inclusive ($8 \le E_T < 36$~GeV) mean isolation energies for the total
cone and its three sub-components (EMCal, inner HCal, outer HCal) at $R=0.3$
and $R=0.4$. All values in GeV.  Note that
\texttt{cluster\_iso\_03\_emcal} is the full EMCal cone integral minus the bare
cluster $E_T$ --- it is a diagnostic, not a decomposition of the total.
% source: isolation_findings.md Sec. "Inclusive shifts (converted - unconverted)"
}
\label{tab:iso_decomposition}
\begin{tabular}{lrrr}
\toprule
Variable & Unconverted & Converted & Converted $-$ Unconverted \\
\midrule
\texttt{cluster\_iso\_03} (total)     & 0.816 & 1.386 & $+0.570$ \\
\texttt{cluster\_iso\_03\_emcal}      & 4.683 & 5.151 & $+0.468$ \\
\texttt{cluster\_iso\_03\_hcalin}     & 0.269 & 0.213 & $-0.056$ \\
\texttt{cluster\_iso\_03\_hcalout}    & 0.287 & 0.253 & $-0.033$ \\
\midrule
\texttt{cluster\_iso\_04} (total)     & 0.957 & 1.531 & $+0.574$ \\
\texttt{cluster\_iso\_04\_emcal}      & 8.166 & 8.484 & $+0.318$ \\
\texttt{cluster\_iso\_04\_hcalin}     & 0.297 & 0.241 & $-0.056$ \\
\texttt{cluster\_iso\_04\_hcalout}    & 0.434 & 0.398 & $-0.036$ \\
\bottomrule
\end{tabular}
\end{table}

The $+0.57$~GeV mean shift at $R=0.3$ is driven by the EMCal sub-component
($+0.47$~GeV) while both HCal sub-components move in the opposite direction
($-0.05$ and $-0.03$~GeV). The physical picture is exactly the converse of
what is ``leaking out'' of the cluster in Sec.~\ref{sec:response}: the
\epem\ arm that curls outside the cluster footprint still deposits its
shower \emph{inside} the $R=0.3$ isolation cone, inflating
\texttt{iso\_03\_emcal}. Meanwhile, because the converted EMCal shower is
more compactly absorbed (two overlapping \epem\ showers), \emph{less} energy
leaks to the HCal. The same physics explains why the effect is largest at
low $E_T$ (where the bent-pair opening is largest relative to a fixed $R=0.3$
cone) and fades at high $E_T$ where the pair is collimated enough to fall
inside the cluster.
% source: isolation_findings.md Sec. "Conclusion -- the shift is driven by the EMCal"


%% ===============================================================
\section{Shower Shape and ID-Variable Separation}
\label{sec:showershape}

The shower-shape BDT uses the 25 features specified in
\texttt{FunWithxgboost/config.yaml}. Here we examine the response of the most
informative features (\wphicog, \wetacog, $e_{32}/e_{35}$, BDT score, CNN
prob) to the converted/unconverted split.

\subsection{Headline: \wphicog\ broadens by $+91\%$}

\begin{figure}[H]
\centering
\begin{subfigure}[t]{0.48\textwidth}
\includegraphics[width=\textwidth]{converted_photon_study/figures/showershape_wphi_cogx_inclusive.pdf}
\caption{Inclusive \wphicog\ ($8 \le E_T < 36$~GeV)}
\end{subfigure}
\hfill
\begin{subfigure}[t]{0.48\textwidth}
\includegraphics[width=\textwidth]{converted_photon_study/figures/showershape_wphi_cogx_mean.pdf}
\caption{Mean \wphicog\ vs.\ \(p_T\)}
\end{subfigure}
\caption{\wphicog\ (phi cluster width) inclusive distribution and
\(p_T\)-binned mean.  Converted photons show a clear bimodal structure ---
one population overlapping the unconverted core (pair unopened) and one at
high \wphicog\ (pair opened). The resulting inclusive mean shift is
$+0.116$ ($+91\%$ relative, $2.15\,\sigma_{\mathrm{unconv}}$).
% source: showershape_findings.md Sec. 3 inclusive table; Sec. 5 ranking; Sec. 6 point 5
}
\label{fig:wphi_cogx_panels}
\end{figure}

The bimodality in Fig.~\ref{fig:wphi_cogx_panels} matches exactly the
bimodal $\Delta\phi$ seen in Fig.~\ref{fig:response_angular}. Both arise from
the same physics: the $\sim 1.4$~T sPHENIX solenoid opens the \epem\ pair in
$\phi$ but not in $\eta$. The resulting two-lobe cluster is wider in the
bending plane, which is detected by \wphicog. The PPG12 parametric tight cut
(\texttt{tight\_wphi\_cogx\_max}, implemented in
\texttt{FunWithxgboost/BDTinput.C}) therefore preferentially rejects
opened-pair converted photons.

\subsection{Other features}

\begin{figure}[H]
\centering
\begin{subfigure}[t]{0.48\textwidth}
\includegraphics[width=\textwidth]{converted_photon_study/figures/showershape_weta_cogx_inclusive.pdf}
\caption{\wetacog\ inclusive}
\end{subfigure}
\hfill
\begin{subfigure}[t]{0.48\textwidth}
\includegraphics[width=\textwidth]{converted_photon_study/figures/showershape_bdt_score_nosplit_inclusive.pdf}
\caption{Shower-shape BDT score inclusive}
\end{subfigure}
\caption{Eta cluster width (\wetacog) and shower-shape BDT score (from
\texttt{CLUSTERINFO\_CEMC\_NO\_SPLIT}). The BDT distribution for converted
photons has a clear low-score tail ($\sim 0.3$--$0.5$) that is essentially
absent for unconverted photons.
% source: showershape_findings.md Sec. 3, Sec. 6 point 2
}
\label{fig:wetacog_bdt}
\end{figure}

\begin{figure}[H]
\centering
\includegraphics[width=0.55\textwidth]{converted_photon_study/figures/showershape_CNN_prob_inclusive.pdf}
\caption{Inclusive CNN photon probability \texttt{CNN\_prob}. Converted
photons are pushed strongly toward zero (mean $0.638 \to 0.464$,
$-27.3\%$ relative). The CNN is not used in the nominal selection but is by
far the most discriminating single variable for converted/unconverted
topology; it is shown here as a cross-check.
% source: showershape_findings.md Sec. 3 and Sec. 6 point 4
}
\label{fig:cnn_prob}
\end{figure}

\begin{table}[H]
\centering
\caption{Inclusive ($8 \le E_T < 36$~GeV) means of selected shower-shape
observables and signed shift (converted $-$ unconverted). The last column
is the absolute shift expressed in units of $\sigma_{\mathrm{unconv}}$,
ranking separation power.
% source: showershape_findings.md Sec. 3 and Sec. 5 discriminating-variables table
}
\label{tab:showershape_inclusive}
\begin{tabular}{lrrrr}
\toprule
Variable       & $\langle\,\rangle$ unconv & $\langle\,\rangle$ conv & $\Delta$ (rel.) & $|\Delta|/\sigma_{\mathrm{u}}$ \\
\midrule
\wetacog                 & 0.191 & 0.225 & $+0.034$ ($+17.8\%$) & 0.27 \\
\wphicog                 & 0.128 & 0.244 & $+0.116$ ($+91.2\%$) & 2.15 \\
$e_{11}/e_{33}$          & 0.649 & 0.597 & $-0.052$ ($-8.0\%$)  & 0.31 \\
$e_{32}/e_{35}$          & 0.980 & 0.946 & $-0.034$ ($-3.4\%$)  & 2.33 \\
$\mathrm{ET1}$ (leading) & 0.939 & 0.911 & $-0.028$ ($-3.0\%$)  & 0.55 \\
$\mathrm{ET3}$           & 0.661 & 0.586 & $-0.075$ ($-11.3\%$) & 0.31 \\
CNN prob                 & 0.638 & 0.464 & $-0.174$ ($-27.3\%$) & 0.40 \\
\texttt{cluster\_prob}   & 0.565 & 0.415 & $-0.149$ ($-26.5\%$) & 0.48 \\
\texttt{merged\_prob}    & 0.804 & 0.673 & $-0.131$ ($-16.3\%$) & 0.55 \\
BDT score (no-split)     & 0.736 & 0.616 & $-0.121$ ($-16.4\%$) & 0.74 \\
\bottomrule
\end{tabular}
\end{table}

The top two discriminators ($e_{32}/e_{35}$ at $2.33\,\sigma$ and \wphicog\
at $2.15\,\sigma$) both test for shower broadening in the bending plane:
$e_{32}/e_{35}$ is the ratio of the central $3\times 2$ to $3\times 5$ tower
energy, which decreases sharply when the \epem\ pair deposits energy in the
outer columns. The BDT score, by contrast, has a separation of only
$0.74\,\sigma$ on the mean but shows a well-separated \emph{low-score tail}
for converted photons (Fig.~\ref{fig:wetacog_bdt}), which drives its
contribution to the tight-selection efficiency.
% source: showershape_findings.md Sec. 5 and Sec. 6

\subsection{BDT degradation vs.\ \(p_T\)}

\begin{table}[H]
\centering
\caption{Shower-shape BDT mean shift $\Delta\langle\mathrm{BDT}\rangle =
\langle\mathrm{BDT}\rangle_{\mathrm{conv}} -
\langle\mathrm{BDT}\rangle_{\mathrm{unconv}}$ vs.\ bin-centre \(p_T\).
The BDT degradation \emph{grows} with \(p_T\), driven by the growing
low-score tail of converted photons. The nominal parametric tight cut is
$\mathrm{BDT} > 0.80 - 0.015 \cdot E_T$
(\texttt{bdt\_min\_intercept}, \texttt{bdt\_min\_slope} in
\texttt{efficiencytool/config\_bdt\_nom.yaml}, lines 154--156).
% source: showershape_findings.md Sec. 4 (pT dependence table); config_bdt_nom.yaml
}
\label{tab:bdt_vs_pt}
\begin{tabular}{rrr}
\toprule
\(p_T\) centre [GeV] & $\Delta\langle\mathrm{BDT}\rangle$ (abs.) & rel.\ (\%) \\
\midrule
 9  & $-0.135$ & $-18.4$ \\
11  & $-0.109$ & $-14.7$ \\
13  & $-0.106$ & $-14.4$ \\
15  & $-0.116$ & $-15.9$ \\
17  & $-0.131$ & $-18.1$ \\
19  & $-0.151$ & $-21.1$ \\
21  & $-0.162$ & $-22.7$ \\
23  & $-0.178$ & $-25.3$ \\
25  & $-0.176$ & $-25.0$ \\
27  & $-0.171$ & $-24.7$ \\
30  & $-0.166$ & $-24.2$ \\
34  & $-0.168$ & $-25.8$ \\
\bottomrule
\end{tabular}
\end{table}

The most-used BDT features (\wetacog, \wphicog, $\mathrm{ET1}$,
$e_{11}/e_{33}$, $e_{32}/e_{35}$, see \texttt{FunWithxgboost/config.yaml}
lines 48--73) all shift in the direction that makes converted photons look
more background-like. The BDT training is therefore implicitly rejecting
converted photons even though no truth conversion information enters the
classifier.
% source: showershape_findings.md Sec. 5


%% ===============================================================
\section{Selection Efficiency by Conversion Category}
\label{sec:efficiency}

Sections~\ref{sec:response}--\ref{sec:showershape} characterised the
response-level, isolation-level and BDT-level \emph{distributions}.
This section translates those distribution shifts into \emph{efficiencies}
through the nominal PPG12 parametric cuts, separately for unconverted
versus converted truth photons. The measurement directly answers
recommendation~\ref{sec:next_steps}(2): what is the efficiency of each
selection stage, and how much of the nominal signal efficiency comes from
each conversion category?

\paragraph{Definitions and caveats.}
For every truth photon in the fiducial region ($|\eta^{\gamma}|<0.7$,
\texttt{particle\_truth\_iso\_03} $<4$~GeV, $8 \le p_T^{\gamma} < 36$~GeV) we
record whether it satisfies each of three stages:
%
\begin{itemize}
  \item $\varepsilon_{\mathrm{reco}}$: at least one
  \texttt{CLUSTERINFO\_CEMC} cluster with
  \texttt{cluster\_truthtrkID}$={}$\texttt{particle\_trkid} and
  $E_T^{\mathrm{cluster}}>8$~GeV (i.e.\ reco'd and inside the analysis
  $p_T$ range).
  \item $\varepsilon_{\mathrm{ID}}$: at least one
  \texttt{CLUSTERINFO\_CEMC\_NO\_SPLIT} cluster with the same matching
  requirement \emph{and} \mbox{BDT $> 0.80-0.015\cdot E_T$}. BDT score is
  only stored on the unsplit node for this sample (same caveat as
  Sec.~\ref{sec:showershape}).
  \item $\varepsilon_{\mathrm{iso}}$: the matched CEMC cluster satisfies
  \texttt{cluster\_iso\_03} $<0.502095+0.0433036\cdot E_T$~GeV. Because
  \texttt{cluster\_iso\_topo\_04} is not populated in the input file, we
  use the $R=0.3$ calo-tower iso with the topo-iso-tuned cut intercept
  --- the same proxy used in Sec.~\ref{sec:isolation}, so absolute
  numbers should be read as illustrative of \emph{shape}, not of the
  nominal analysis pass fraction.
  \item $\varepsilon_{\mathrm{total}}$: all three simultaneously (truth
  photon has both a CEMC- and a NO\_SPLIT-matched cluster with both cuts
  satisfied).
\end{itemize}
%
All absolute efficiencies are normalised per fiducial truth photon in
the category (i.e.\ divided by $N_{\mathrm{truth}}$ for that category).
Conditional efficiencies $\varepsilon_{\mathrm{ID}\mid\mathrm{reco}}$
and $\varepsilon_{\mathrm{iso}\mid\mathrm{reco}}$ use the
same-node reco denominator to avoid the CEMC/NO\_SPLIT cluster-count
mismatch.

\paragraph{Per-stage efficiency vs.\ $p_T^{\gamma}$.}
Figure~\ref{fig:eff_abs} shows the four absolute efficiencies as a
function of truth $p_T$ for unconverted (black) and converted (red) photons.
Figure~\ref{fig:eff_cond} gives the conditional forms
$\varepsilon_{\mathrm{ID}\mid\mathrm{reco}}$ and
$\varepsilon_{\mathrm{iso}\mid\mathrm{reco}}$.
Figure~\ref{fig:eff_ratio} shows the per-stage ratio
$\varepsilon_{\mathrm{conv}}/\varepsilon_{\mathrm{unconv}}$.

\begin{figure}[H]
\centering
\includegraphics[width=0.95\textwidth]{converted_photon_study/figures/efficiency_vs_pt.pdf}
\caption{Absolute selection efficiency versus truth $p_T$ per conversion
category, four stages. Denominator is the fiducial truth-photon count
within each category. The drop from left to right shows the sequential
loss from reco $\to$ BDT $\to$ iso. Converted photons consistently
track below the unconverted curves.}
\label{fig:eff_abs}
\end{figure}

\begin{figure}[H]
\centering
\includegraphics[width=0.95\textwidth]{converted_photon_study/figures/efficiency_cond_vs_pt.pdf}
\caption{Conditional efficiencies:
$\varepsilon_{\mathrm{ID}\mid\mathrm{reco}}$ (left, denominator =
NO\_SPLIT-matched) and $\varepsilon_{\mathrm{iso}\mid\mathrm{reco}}$
(right, denominator = CEMC-matched). Converted photons lose a larger
fraction of the surviving sample at both the BDT cut and the iso cut.}
\label{fig:eff_cond}
\end{figure}

\begin{figure}[H]
\centering
\includegraphics[width=0.65\textwidth]{converted_photon_study/figures/efficiency_ratio_vs_pt.pdf}
\caption{Efficiency ratio
$\varepsilon_{\mathrm{conv}}/\varepsilon_{\mathrm{unconv}}$ per stage.
Values below unity quantify the conversion penalty at each stage.}
\label{fig:eff_ratio}
\end{figure}

\paragraph{Inclusive summary.}
Table~\ref{tab:eff_inclusive} lists the inclusive ($8\le p_T<36$~GeV)
efficiencies per category along with the conv/unconv ratio.

\begin{table}[H]
\centering
\caption{Inclusive ($8 \le p_T^{\gamma}<36$~GeV) selection efficiency per
conversion category. Numerator definitions are given in the text;
denominator is $N_{\mathrm{truth}}$ for that category (or the
same-node reco match for conditional rows).
% source: rootFiles/efficiencies_summary.json (analyze_efficiencies.py)
}
\label{tab:eff_inclusive}
\begin{tabular}{lrrr}
\toprule
Stage & Unconverted & Converted & Ratio (conv/unconv) \\
\midrule
$N_{\mathrm{truth}}$                                     & 803,718 & 168,685 & --- \\
$\varepsilon_{\mathrm{reco}}$ (CEMC match, $E_T>8$ GeV)  & 64.3\% & 51.9\% & 0.808 \\
$\varepsilon_{\mathrm{ID}}$ (NO\_SPLIT match, BDT pass)  & 50.7\% & 29.6\% & 0.584 \\
$\varepsilon_{\mathrm{iso}}$ (CEMC iso pass)             & 41.7\% & 26.1\% & 0.626 \\
$\varepsilon_{\mathrm{total}}$ (reco $\wedge$ BDT $\wedge$ iso) & 32.5\% & 15.5\% & 0.477 \\
$\varepsilon_{\mathrm{ID}\mid\mathrm{reco}}$             & 78.8\% & 57.0\% & 0.723 \\
$\varepsilon_{\mathrm{iso}\mid\mathrm{reco}}$            & 64.9\% & 50.2\% & 0.775 \\
\bottomrule
\end{tabular}
\end{table}

\paragraph{Physics interpretation.}
Three observations follow from the numbers above:
%
\begin{enumerate}
  \item The conversion penalty at the reco stage is modest ---
  the ratio $\varepsilon_{\mathrm{reco}}^{\mathrm{conv}}/
  \varepsilon_{\mathrm{reco}}^{\mathrm{unconv}}\approx 0.81$.
  Part of this loss is the response-tail of Sec.~\ref{sec:response}
  migrating matched clusters below the 8~GeV threshold.
  \item The BDT cut is where the largest loss happens: the ratio drops
  to $0.58$. This is consistent with the $-16\%$ inclusive
  BDT mean shift and the growing low-score tail (Table~\ref{tab:bdt_vs_pt}).
  \item The iso cut contributes an additional conditional factor
  $0.77$ (given a reco'd cluster). Combined with the BDT
  penalty, the total efficiency ratio is $0.48$ --- converted
  photons are roughly half as efficient as unconverted ones end-to-end.
\end{enumerate}
%
The nominal signal efficiency already integrates over this mix at the
natural $17\%$ converted fraction, so the nominal result is
self-consistent provided the GEANT material budget matches reality. The
table supplies the quantity that drives the material-budget systematic
numerically: $\Delta\varepsilon_{\mathrm{total}} =
(\varepsilon_{\mathrm{total}}^{\mathrm{conv}}-
\varepsilon_{\mathrm{total}}^{\mathrm{unconv}})\cdot\Delta f_{\mathrm{conv}}$
per unit shift in the simulated conversion fraction.


%% ===============================================================
\section{Physics Interpretation}
\label{sec:physics}

The four analyses (fractions, response, isolation, shower shape) assemble
into a single, self-consistent picture of converted photons in sPHENIX
dictated by the $\sim 1.4$~T solenoid:

\begin{enumerate}
  \item The conversion probability is flat in \(p_T\) and rises with
  $|\eta|$. The $\sim 17\%$ central-barrel value is consistent with the
  quoted $10$--$20\%\,X_0$ of integrated material upstream of the EMCal.
  % source: fractions_findings.md Sec. "Per-|eta|-bin" and Sec. "Physics discussion"
  \item The magnetic field opens the \epem\ pair in $\phi$ (but not in
  $\eta$). One of the two arms has a finite probability of curling outside
  the cluster footprint, leaving partial energy in the EMCal. This produces
  a long low-$R$ tail in the energy response --- the Gaussian core shifts by
  only $+0.003$ while the Crystal Ball $\mu$ moves by $-0.028$ and the CB
  $\sigma$ broadens by $+35\%$, consistent with a sub-peak below unity
  coexisting with a two-arm core still close to $R=1$ --- together with a
  bimodal $\Delta\phi$ distribution for the converted category.
  % source: Table 2 (Gaussian + Crystal Ball fit parameters); Sec. 4.3
  \item The lost-arm shower lands inside the $R=0.3$ isolation cone and
  inflates \texttt{cluster\_iso\_03\_emcal} by $\sim 0.5$~GeV on average.
  Meanwhile the HCal sub-components move in the opposite direction (the
  two overlapping \epem\ showers are absorbed more compactly in the EMCal),
  confirming the same spatial leakage pattern. The parametric iso cut
  therefore rejects converted photons $\sim 14.6$~pp more often than
  unconverted ones inclusively, rising to $\sim 23$~pp at low \(p_T\).
  % source: isolation_findings.md Sec. "Conclusion -- the shift is driven by the EMCal"
  \item The wider, two-lobed transverse shower profile is picked up most
  cleanly by \wphicog\ (mean bimodality, $+91\%$ relative shift) and
  $e_{32}/e_{35}$ (tightens the central $3\times 2$ fraction). The BDT
  absorbs these features through its training on jet-MC background and
  produces a well-separated low-score tail whose effect on the tight
  selection grows with \(p_T\), reaching $-26\%$ at 23--34~GeV.
  % source: showershape_findings.md Sec. 4 and Sec. 5
  \item All four conversion signatures share a common \(p_T\) evolution:
  the effect is strongest at low \(p_T\) (large magnetic opening relative to
  the cluster/cone scale) and fades at high \(p_T\) where the \epem\ pair is
  collimated. The BDT shift is the exception --- it grows with \(p_T\) ---
  because the parametric cut itself becomes tighter at high \(p_T\) and the
  low-score tail of converted photons sits on the wrong side of an
  increasingly stringent boundary.
\end{enumerate}


%% ===============================================================
\section{Implications for the Isolated-Photon Cross-Section}
\label{sec:cross_section}

The converted-photon signatures above enter the PPG12 cross-section
measurement as an ID$+$iso efficiency effect:

\begin{itemize}
  \item The current nominal efficiency $\varepsilon(\mathrm{reco}\to
  \mathrm{tight+iso})$ is computed on simulated signal events that
  \emph{already} contain conversions at their natural $\sim 17\%$ rate. As
  long as the GEANT material budget matches reality, the effect is
  correctly absorbed into the nominal efficiency.
  \item The converted population is simultaneously penalised by the BDT
  tight cut (lower mean BDT, growing low-score tail) \emph{and} by the iso
  parametric cut (EMCal pollution from the lost \epem\ arm). The PPG12
  signal efficiency is therefore a product of two partially anticorrelated
  selection factors, both of which track the converted fraction.
  \item A mismatch between the simulated and data material budget (e.g.
  missing support structure, cable run, coolant volume) would produce a
  \emph{double} shift: the effective converted fraction would move, and
  both $\varepsilon_{\mathrm{BDT}}$ and $\varepsilon_{\mathrm{iso}}$ would
  respond.
\end{itemize}

\noindent The conclusion is that the material budget is one of the key
cross-section systematics to control.  A dedicated material-budget
uncertainty --- propagated through the converted fraction to both
$\varepsilon_{\mathrm{BDT}}$ and $\varepsilon_{\mathrm{iso}}$ --- would
correctly cover the combined ID$+$iso efficiency response.


%% ===============================================================
\section{Recommendations and Next Steps}
\label{sec:next_steps}

\begin{enumerate}
  \item \textbf{Re-run the isolation study with topo-iso branches populated.}
  The nominal analysis uses \texttt{cluster\_iso\_topo\_04}
  (\texttt{use\_topo\_iso: 2}); the current study used \texttt{cluster\_iso\_03}
  by necessity. Populating \texttt{cluster\_iso\_topo\_03/04} on this sample
  (re-running \texttt{apply\_BDT.C} or the slim-tree builder with the topo
  branches enabled) will allow the absolute pass-fraction numbers to be
  tied directly to the nominal cut.
  % source: isolation_findings.md Sec. "Concerns and next steps"
  \item \textbf{Split the tight-selection efficiency by conversion category.}
  Apply $\mathrm{BDT} > 0.80 - 0.015 \cdot E_T$ and the parametric iso cut
  of Eq.~\ref{eq:iso_cut} separately to unconverted, converted, and
  combined populations to obtain
  $\varepsilon_{\mathrm{unconv}}$, $\varepsilon_{\mathrm{conv}}$, and the
  natural-fraction-weighted efficiency. The difference between the
  weighted and the unconverted-only result is the conversion contribution
  to the nominal efficiency.
  % source: showershape_findings.md Sec. 8
  \item \textbf{Cross-check BDT data/MC agreement at low BDT values.} The
  converted-photon population produces a distinct low-score BDT tail. The
  shape of that tail in \emph{data} (which carries no truth conversion
  flag) is a direct probe of the simulated material budget: if the MC
  underestimates the material, the observed low-BDT-score rate will be
  \emph{higher} than predicted.
  \item \textbf{Independent toy: artificially remove conversions from MC
  and propagate.} A zero-conversion counterfactual (e.g.\ selecting only
  \texttt{particle\_converted==0}) gives a direct handle on the converted
  contribution to the nominal efficiency and to the cross-section.
  \item \textbf{Repeat on photon5 and photon20.} The pT dependence of the
  converted-efficiency loss should be cross-checked at 1--14~GeV
  (\texttt{photon5}) and 20--40~GeV (\texttt{photon20}) to confirm the
  crossover at $p_T \sim 25$~GeV and the growing BDT shift.
  % source: showershape_findings.md Sec. 8
\end{enumerate}


%% ===============================================================
\appendix
\section{Slimtree branches used}
\label{app:branches}

\begin{itemize}
  \item \texttt{particle\_pid}, \texttt{particle\_Pt}, \texttt{particle\_Eta},
    \texttt{particle\_Phi}, \texttt{particle\_E}, \texttt{particle\_trkid}
    --- truth four-vector and track id (written at
    \texttt{CaloAna24.cc} $\sim$957--966).
  \item \texttt{particle\_converted} --- conversion flag
    ($0$/$1$/$2$; see Sec.~\ref{sec:intro}).
    % source: CaloAna24.cc lines 953-959, 970
  \item \texttt{particle\_truth\_iso\_03} --- truth isolation $E_T$ in $R=0.3$.
  \item \texttt{CLUSTERINFO\_CEMC.cluster\_Et},
    \texttt{cluster\_truthtrkID} --- reconstructed cluster $E_T$ and
    truth-track assignment.
  \item Shower shape: \texttt{weta\_cogx}, \texttt{wphi\_cogx},
    \texttt{et1}--\texttt{et4}, \texttt{e11/e33}, \texttt{e32/e35},
    \texttt{CNN\_prob}, \texttt{cluster\_prob}, \texttt{merged\_prob}.
  \item \texttt{CLUSTERINFO\_CEMC\_NO\_SPLIT.bdt\_score} --- shower-shape BDT
    output.
  \item Isolation: \texttt{cluster\_iso\_03}, \texttt{cluster\_iso\_04},
    \texttt{..\_emcal}, \texttt{..\_hcalin}, \texttt{..\_hcalout}.
\end{itemize}

\section{Figure inventory}
\label{app:figure_inventory}

All figures live under
\texttt{reports/converted\_photon\_study/figures/}. A compact inventory:

\begin{itemize}
  \item Fractions: \texttt{\seqsplit{fractions\_vs\_pt\_eta07.pdf}},
    \texttt{\seqsplit{fractions\_vs\_abseta\_eta07.pdf}},
    \texttt{\seqsplit{fractions\_match\_eff\_eta07.pdf}} (plus
    \texttt{\_eta15} analogues; \texttt{fractions\_2d\_*};
    \texttt{\seqsplit{fractions\_multiplicity\_*}}).
  \item Response: \texttt{\seqsplit{response\_shape\_grid.pdf}},
    \texttt{\seqsplit{response\_shape\_ET\_grid.pdf}},
    \texttt{\seqsplit{response\_mean\_vs\_pt.pdf}},
    \texttt{\seqsplit{response\_sigma\_vs\_pt.pdf}},
    \texttt{\seqsplit{response\_cb\_params\_vs\_pt.pdf}},
    \texttt{\seqsplit{response\_tailfrac\_vs\_pt.pdf}},
    \texttt{\seqsplit{response\_match\_eff\_vs\_pt.pdf}},
    \texttt{\seqsplit{response\_eta\_dep.pdf}},
    \texttt{\seqsplit{response\_angular\_deltas.pdf}}.
  \item Isolation: \texttt{\seqsplit{isolation\_cluster\_iso\_03.pdf}},
    \texttt{\seqsplit{isolation\_cluster\_iso\_03\_emcal.pdf}} (plus
    \texttt{\_hcalin}, \texttt{\_hcalout}, \texttt{\_cumulative}; $R=0.4$
    analogues), \texttt{\seqsplit{isolation\_mean\_vs\_pt\_R003.pdf}},
    \texttt{\_R004.pdf},
    \texttt{\seqsplit{isolation\_pass\_fraction\_vs\_pt.pdf}}.
  \item Shower shape (each variable has \texttt{\_inclusive}, \texttt{\_mean}
    and a 12-panel $p_T$-binned PDF):
    \texttt{\seqsplit{showershape\_weta\_cogx*.pdf}},
    \texttt{\seqsplit{showershape\_wphi\_cogx*.pdf}},
    \texttt{showershape\_et1*.pdf} to \texttt{et4*.pdf},
    \texttt{\seqsplit{showershape\_e11\_over\_e33*.pdf}},
    \texttt{\seqsplit{showershape\_e32\_over\_e35*.pdf}},
    \texttt{\seqsplit{showershape\_et2\_over\_et1*.pdf}} to
    \texttt{et4\_over\_et1*.pdf},
    \texttt{\seqsplit{showershape\_CNN\_prob*.pdf}},
    \texttt{\seqsplit{showershape\_cluster\_prob*.pdf}},
    \texttt{\seqsplit{showershape\_merged\_prob*.pdf}},
    \texttt{\seqsplit{showershape\_bdt\_score\_nosplit*.pdf}}.
\end{itemize}

\end{document}
